\chapter{Resultados y discusión}
\label{cap:resultados}

Este capítulo presenta los resultados finales de la fase experimental definida en el Capítulo~\ref{cap:experimentos}. El análisis usa el subconjunto balanceado del archivo \textttbreak{benchmark\_final.csv}: datasets 1--5, semillas 42, 84 y 126, y sólo modelos con las 15 corridas completas. En total, el subconjunto final contiene 360 corridas, correspondientes a 24 modelos.

\hfill

Las tablas reportan la media y el intervalo de confianza de la distribución t de Student al 95\% sobre 15 corridas individuales (5 datasets $\times$ 3 semillas), tratándolas como observaciones independientes para aumentar la potencia estadística de los análisis. Valores menores indican mejor desempeño. La métrica principal es MSE; RMSE y MAE se usan como métricas complementarias.

\section{Cobertura final del benchmark}
\label{sec:cobertura-benchmark-final}

La Tabla~\ref{tab:cobertura-final} resume la cobertura utilizada para el análisis final. El archivo de resultados contiene algunas filas adicionales de corridas parciales, pero no se usan para las conclusiones porque no forman parte del conjunto balanceado.

\begin{table}[H]
	\centering
	\caption{Cobertura final usada en el análisis.}
	\label{tab:cobertura-final}
	\begin{tabular}{lcc}
		\toprule
		Grupo & Modelos completos & Corridas completas \\
		\midrule
		Variantes \CTTransformer\ encoder & 10 & 150 \\
		Variantes encoder--decoder optimizadas & 8 & 120 \\
		Baselines de estado del arte adaptados & 2 & 30 \\
		Baselines simples & 2 & 30 \\
		Ablaciones & 2 & 30 \\
		\midrule
		Total & 24 & 360 \\
		\bottomrule
	\end{tabular}
\end{table}

\section{Metodología de presentación y análisis de resultados}
\label{sec:metodologia-analisis-resultados}

Para asegurar la validez científica y el rigor de las comparaciones presentadas en este capítulo, se adoptan las siguientes definiciones y metodologías estadísticas:

\begin{itemize}
	\item \textbf{Intervalos de Confianza t-Student}: En lugar del formato tradicional de media más/menos desviación estándar (que solo describe la dispersión muestral), cada métrica principal se reporta mediante su media y su intervalo de confianza al 95\% para la media poblacional. Este se calcula utilizando la distribución t de Student con $n-1$ grados de libertad:
	$$\bar{x} \pm t_{\alpha/2, n-1} \frac{s}{\sqrt{n}}$$
	donde $\bar{x}$ es la media muestral, $s$ es la desviación estándar muestral, $n = 15$ es el número de corridas independientes (5 datasets $\times$ 3 semillas) y $t_{\alpha/2, 14} \approx 2.145$ es el valor crítico para $\alpha = 0.05$ bilateral. Este intervalo define el rango en el cual se sitúa el verdadero error medio esperado del modelo con un 95\% de certeza.
	
	\item \textbf{Pruebas de Significancia Estadística Emparejadas}: Dado que todos los modelos se evalúan exactamente sobre el mismo conjunto de datasets y semillas, los resultados de desempeño están emparejados. Por ello, para contrastar si la superioridad del modelo ganador es estadísticamente significativa frente a las alternativas, se aplican dos pruebas bilaterales:
	\begin{enumerate}
		\item \emph{Prueba t de Student para muestras emparejadas}: Evalúa si la diferencia media de los errores entre dos modelos es significativamente distinta de cero, asumiendo normalidad en las diferencias de error.
		\item \emph{Prueba de rangos con signo de Wilcoxon}: Alternativa no paramétrica que no asume normalidad, basada en el orden de los rangos de las diferencias.
	\end{enumerate}
	Un p-value ($p$) menor a $0.05$ se considera suficiente para rechazar la hipótesis nula ($H_0$) de igualdad de rendimiento, declarando una ventaja estadísticamente significativa.
	
	\item \textbf{Métricas de Error}: Se reportan tres métricas complementarias para evaluar distintas propiedades de las predicciones:
	\begin{enumerate}
		\item \emph{MSE (Mean Squared Error)}: Métrica principal de optimización. Al elevar los errores al cuadrado, penaliza de forma muy estricta las desviaciones grandes, lo cual es crítico en aplicaciones reales donde los errores extremos son costosos.
		\item \emph{RMSE (Root Mean Squared Error)}: Mantiene las propiedades del MSE pero escala el error a la misma unidad física de la variable objetivo.
		\item \emph{MAE (Mean Absolute Error)}: Promedio de los errores absolutos. Es más robusto frente a valores atípicos y describe el error promedio típico.
	\end{enumerate}
	
	\item \textbf{Victorias}: Indica la proporción de datasets (en este caso, de los 5 datasets balanceados) en los cuales el modelo respectivo obtuvo el menor MSE promedio.
\end{itemize}

\section{Resultado principal}
\label{sec:resultado-principal-final}

El mejor desempeño observado en la evaluación final corresponde a \textttbreak{EncDec-Opt-Small-MT8\_FT\_AR\_Contiguous}. Esta variante obtiene el menor MSE, RMSE y MAE, y gana en los 5 de 5 datasets al comparar promedios por dataset.

\hfill

Arquitectónicamente, el modelo ganador es un encoder--decoder pequeño: $d_{\text{model}}=64$, 2 capas de encoder, 4 cabezas de atención, \emph{feed-forward} de dimensión 128 y \emph{dropout} 0.1. A diferencia del encoder puro \CTTransformer, procesa la historia en un encoder y los tokens objetivo en un decoder liviano de una capa con atención cruzada hacia la historia. La sigla MT8 indica que se entrenó con ocho tokens objetivo. El sufijo \textttbreak{FT\_AR\_Contiguous} indica un ajuste posterior autoregresivo con máscara causal y objetivos futuros contiguos.

\hfill

La Tabla~\ref{tab:resultados-finales-principales} compara el modelo ganador con las variantes y baselines más relevantes. El tiempo reportado para modelos con ajuste autoregresivo incluye tanto el entrenamiento base como el ajuste posterior.

\begin{table}[H]
	\centering
	\caption{Resultados finales de modelos representativos.}
	\label{tab:resultados-finales-principales}
	\scriptsize
	\resizebox{\textwidth}{!}{%
	\begin{tabular}{lcccccc}
		\toprule
		Modelo & MSE & RMSE & MAE & Parámetros & Tiempo medio & Victorias \\
		\midrule
		EncDec-Opt pequeño MT8 + AR contiguo & \textbf{0.1979 [0.1702, 0.2257]} & \textbf{0.4417 [0.4110, 0.4724]} & \textbf{0.2193 [0.1949, 0.2437]} & 117.8K & 0.46 h & 5/5 \\
		\CTTransformer\ pequeño + Time2Vec + sesgo temporal & 0.2807 [0.2443, 0.3170] & 0.5265 [0.4927, 0.5603] & 0.2906 [0.2519, 0.3293] & 67.8K & 0.74 h & 0/5 \\
		\CTTransformerSmall & 0.2817 [0.2456, 0.3178] & 0.5275 [0.4940, 0.5611] & 0.2935 [0.2526, 0.3343] & 67.4K & 0.69 h & 0/5 \\
		\CTTransformerBase & 0.2906 [0.2568, 0.3244] & 0.5363 [0.5054, 0.5673] & 0.2859 [0.2471, 0.3247] & 636.2K & 0.92 h & 0/5 \\
		Modelo lineal temporal & 0.2919 [0.2558, 0.3281] & 0.5372 [0.5043, 0.5702] & 0.3286 [0.2871, 0.3700] & 10.8K & 0.69 h & 0/5 \\
		STraTS adaptado & 0.3274 [0.2765, 0.3783] & 0.5668 [0.5219, 0.6118] & 0.3205 [0.2763, 0.3647] & 229.2K & 0.63 h & 0/5 \\
		Persistencia & 0.4201 [0.3144, 0.5258] & 0.6365 [0.5664, 0.7066] & 0.3435 [0.2530, 0.4341] & 0 efectivos & 0.00 h & 0/5 \\
		CoFormer univariado & 0.5663 [0.5198, 0.6127] & 0.7506 [0.7200, 0.7812] & 0.4464 [0.3883, 0.5046] & 693.4K & 0.22 h & 0/5 \\
		\bottomrule
	\end{tabular}}
\end{table}

La reducción de MSE del modelo ganador es de 29.7\% frente a \CTTransformerSmall, 31.9\% frente a \CTTransformerBase, 32.2\% frente al modelo lineal temporal, 39.5\% frente a STraTS adaptado y 52.9\% frente a persistencia. Además, la ventaja no se concentra en un solo dataset: el modelo ganador obtiene el menor MSE promedio en cada uno de los cinco datasets.

\section{Comparación contra baselines}
\label{sec:comparacion-baselines-final}

La comparación contra persistencia muestra que todos los modelos neuronales principales aprenden más que repetir el último valor observado. El resultado más fuerte lo entrega el encoder--decoder autoregresivo: reduce el MSE medio desde 0.4201 hasta 0.1979.

\hfill

La comparación contra el modelo lineal temporal es más informativa. El lineal alcanza MSE 0.2919, muy cerca de \CTTransformerBase\ y de algunas variantes encoder puras. Esto confirma que el benchmark no es trivial: una parametrización simple con historia reciente y tiempos relativos ya captura parte importante de la dinámica. Aun así, el modelo ganador queda claramente por debajo de ese baseline.

\hfill

STraTS adaptado obtiene MSE 0.3274, competitivo pero por debajo de las mejores variantes \CTTransformer. CoFormer univariado queda rezagado con MSE 0.5663. Esta diferencia debe interpretarse con cuidado: CoFormer fue diseñado para explotar relaciones entre múltiples variables asincrónicas, mientras que la evaluación final consolidada es univariada. En este contexto pierde buena parte de su sesgo inductivo original.

\section{Variantes encoder--decoder}
\label{sec:variantes-encoder-decoder-final}

La Tabla~\ref{tab:encdec-final} muestra que el ajuste autoregresivo es el cambio arquitectónico más decisivo de la evaluación. La variante encoder--decoder pequeña sin ajuste autoregresivo obtiene MSE 0.2945, similar a \CTTransformerBase. Después del ajuste autoregresivo contiguo, el MSE baja a 0.1979, una reducción aproximada de 33\%.

\begin{table}[H]
	\centering
	\caption{Resultados finales de variantes encoder--decoder optimizadas.}
	\label{tab:encdec-final}
	\scriptsize
	\resizebox{\textwidth}{!}{%
	\begin{tabular}{lccccc}
		\toprule
		Modelo & MSE & RMSE & MAE & Parámetros & Tiempo medio \\
		\midrule
		EncDec-Opt-Small-MT8 & 0.2945 [0.2575, 0.3314] & 0.5395 [0.5060, 0.5730] & 0.2859 [0.2491, 0.3226] & 117.8K & 0.30 h \\
		EncDec-Opt-Small-MT8 + AR contiguo & \textbf{0.1979 [0.1702, 0.2257]} & \textbf{0.4417 [0.4110, 0.4724]} & \textbf{0.2193 [0.1949, 0.2437]} & 117.8K & 0.46 h \\
		EncDec-Opt-Small-MT8 + AR mixto & 0.2043 [0.1767, 0.2320] & 0.4489 [0.4183, 0.4795] & 0.2248 [0.1994, 0.2501] & 117.8K & 0.47 h \\
		EncDec-Opt-Small-MT8 + AR aleatorio & 0.2206 [0.1885, 0.2527] & 0.4659 [0.4317, 0.5001] & 0.2369 [0.2075, 0.2664] & 117.8K & 0.47 h \\
		EncDec-Opt-MT8 & 0.3028 [0.2642, 0.3413] & 0.5470 [0.5125, 0.5814] & 0.2909 [0.2528, 0.3289] & 1.06M & 0.53 h \\
		EncDec-Opt-MT8 + AR contiguo & 0.2102 [0.1775, 0.2429] & 0.4544 [0.4195, 0.4893] & 0.2240 [0.1978, 0.2502] & 1.06M & 0.76 h \\
		EncDec-Opt-MT8 + AR mixto & 0.2172 [0.1833, 0.2512] & 0.4619 [0.4260, 0.4977] & 0.2285 [0.1997, 0.2572] & 1.06M & 0.79 h \\
		EncDec-Opt-MT8 + AR aleatorio & 0.2424 [0.1993, 0.2855] & 0.4866 [0.4437, 0.5295] & 0.2430 [0.2095, 0.2764] & 1.06M & 0.79 h \\
		\bottomrule
	\end{tabular}}
\end{table}

El patrón es consistente: las versiones autoregresivas superan a sus respectivas bases, el modo contiguo es el mejor entre los tres modos de ajuste y el tamaño pequeño domina al tamaño base. Esto sugiere que, para estos cinco datasets univariados, el beneficio principal no proviene de aumentar capacidad, sino de alinear el entrenamiento con la generación secuencial de horizontes futuros.

\hfill

Hay una cautela metodológica importante. Las variantes autoregresivas se evalúan con generación sobre ocho objetivos contiguos, mientras que los modelos no autoregresivos se evalúan con hasta cuatro objetivos dentro del rango de offsets configurado. Por ello, el resultado permite declarar al encoder--decoder pequeño con AR contiguo como el mejor sistema observado en esta evaluación, pero no aísla completamente si la mejora proviene del decoder, del ajuste autoregresivo, del protocolo de objetivos o de la combinación de estos factores.

\section{Variantes \CTTransformer}
\label{sec:variantes-ct-final}

Dentro de las variantes encoder puras de \CTTransformer, la señal principal es que los modelos pequeños son suficientes y, en MSE, levemente mejores que las versiones base. La Tabla~\ref{tab:variantes-ct-final} resume las variantes más relevantes.

\begin{table}[H]
	\centering
	\caption{Resultados finales de variantes encoder de \CTTransformer.}
	\label{tab:variantes-ct-final}
	\scriptsize
	\resizebox{\textwidth}{!}{%
	\begin{tabular}{lcccc}
		\toprule
		Modelo & MSE & RMSE & MAE & Parámetros \\
		\midrule
		\CTTransformer\ pequeño + Time2Vec + sesgo temporal & \textbf{0.2807 [0.2443, 0.3170]} & \textbf{0.5265 [0.4927, 0.5603]} & 0.2906 [0.2519, 0.3293] & 67.8K \\
		\CTTransformer\ pequeño + sesgo temporal & 0.2808 [0.2447, 0.3168] & 0.5266 [0.4931, 0.5602] & 0.2915 [0.2529, 0.3300] & 67.6K \\
		\CTTransformer\ pequeño + resumen atencional & 0.2817 [0.2455, 0.3178] & 0.5275 [0.4940, 0.5610] & 0.2934 [0.2541, 0.3326] & 82.1K \\
		\CTTransformerSmall & 0.2817 [0.2456, 0.3178] & 0.5275 [0.4940, 0.5611] & 0.2935 [0.2526, 0.3343] & 67.4K \\
		\CTTransformer\ pequeño + Time2Vec & 0.2824 [0.2463, 0.3184] & 0.5282 [0.4948, 0.5616] & 0.2938 [0.2540, 0.3335] & 67.8K \\
		\CTTransformer\ base + Time2Vec + sesgo temporal & 0.2851 [0.2498, 0.3203] & 0.5309 [0.4984, 0.5634] & \textbf{0.2756 [0.2395, 0.3116]} & 794.9K \\
		\CTTransformer\ base + sesgo temporal & 0.2846 [0.2483, 0.3208] & 0.5302 [0.4969, 0.5636] & 0.2777 [0.2411, 0.3143] & 794.4K \\
		\CTTransformer\ base + Time2Vec & 0.2867 [0.2505, 0.3228] & 0.5323 [0.4990, 0.5655] & 0.2773 [0.2381, 0.3165] & 794.9K \\
		\CTTransformer\ base + resumen atencional & 0.2869 [0.2506, 0.3231] & 0.5324 [0.4991, 0.5658] & 0.2799 [0.2406, 0.3193] & 852.1K \\
		\CTTransformerBase & 0.2906 [0.2568, 0.3244] & 0.5363 [0.5054, 0.5673] & 0.2859 [0.2471, 0.3247] & 636.2K \\
		\bottomrule
	\end{tabular}}
\end{table}

La mejor variante encoder pura por MSE es \CTTransformer\ pequeño con Time2Vec y sesgo temporal, pero su ventaja frente a \CTTransformerSmall\ es sólo de 0.4\% en MSE. Por lo tanto, la conclusión razonable no es que Time2Vec más sesgo temporal sea decisivo, sino que la familia pequeña de \CTTransformer\ es robusta y difícil de distinguir con sólo cinco datasets.

\hfill

En MAE, las variantes base con Time2Vec o sesgo temporal mejoran al modelo pequeño. Esto sugiere que mayor capacidad y representaciones temporales más flexibles pueden reducir errores absolutos típicos, aunque no necesariamente los errores cuadráticos grandes que dominan el MSE.

\section{Ablaciones}
\label{sec:estudios-ablacion-final}

La Tabla~\ref{tab:ablaciones-final} muestra las dos ablaciones principales. La eliminación del identificador historia--objetivo degrada el desempeño: el MSE sube desde 0.2906 hasta 0.3028, una pérdida cercana a 4.0\%. Esta señal respalda la decisión de marcar explícitamente qué tokens son observaciones pasadas y cuáles son consultas futuras.

\begin{table}[H]
	\centering
	\caption{Ablaciones finales respecto de \CTTransformerBase.}
	\label{tab:ablaciones-final}
	\scriptsize
	\resizebox{\textwidth}{!}{%
	\begin{tabular}{lcccc}
		\toprule
		Modelo & MSE & RMSE & MAE & Parámetros \\
		\midrule
		\CTTransformerBase & 0.2906 [0.2568, 0.3244] & 0.5363 [0.5054, 0.5673] & \textbf{0.2859 [0.2471, 0.3247]} & 636.2K \\
		Sin encoding temporal continuo & \textbf{0.2903 [0.2534, 0.3271]} & \textbf{0.5356 [0.5021, 0.5691]} & 0.2957 [0.2579, 0.3335] & 583.4K \\
		Sin identificador de token objetivo & 0.3028 [0.2698, 0.3357] & 0.5478 [0.5177, 0.5778] & 0.3027 [0.2680, 0.3374] & 600.7K \\
		\CTTransformerSmall & 0.2817 [0.2456, 0.3178] & 0.5275 [0.4940, 0.5611] & 0.2935 [0.2526, 0.3343] & 67.4K \\
		\bottomrule
	\end{tabular}}
\end{table}

El efecto del encoding temporal continuo es menos claro. La ablación sin tiempo continuo queda prácticamente empatada con \CTTransformerBase\ en MSE y RMSE, aunque empeora en MAE. En estos cinco datasets univariados, no hay evidencia suficiente para afirmar que el encoding temporal continuo sea por sí solo el componente decisivo. La evidencia sí favorece el uso del identificador de token objetivo.

\hfill

La comparación con \CTTransformerSmall\ también muestra que reducir capacidad mejora el MSE frente al modelo base. Esto sugiere que las series evaluadas no requieren necesariamente un encoder de mayor tamaño, o que la regularización implícita de la versión pequeña resulta favorable.

\section{Costo computacional}
\label{sec:analisis-costo-final}

El resultado de costo es favorable para el modelo ganador. Aunque requiere dos etapas, entrenamiento base y ajuste autoregresivo, su tiempo medio total es 0.46 horas por corrida y usa 117.8K parámetros. Es más costoso que persistencia, pero más barato que \CTTransformerSmall\ en el registro final y bastante más pequeño que \CTTransformerBase.

\hfill

El aumento de tamaño no entrega una mejora clara. \CTTransformerBase\ usa 636.2K parámetros y obtiene MSE 0.2906, mientras que \CTTransformerSmall\ usa 67.4K y obtiene MSE 0.2817. Algo similar ocurre en encoder--decoder: la variante pequeña con ajuste AR contiguo supera a la variante base con el mismo ajuste, pese a tener aproximadamente una novena parte de los parámetros.

\section{Significancia estadística y alcance}
\label{sec:significancia-estadistica-final}

Para robustecer las conclusiones sobre el desempeño relativo de los modelos y mitigar la limitación del tamaño de muestra al promediar semillas, se trataron las 15 corridas (5 datasets $\times$ 3 semillas) como observaciones independientes emparejadas por par (dataset, semilla). Sobre estas 15 corridas independientes se aplicaron dos pruebas estadísticas bilaterales para evaluar la significancia de la diferencia en MSE: la prueba t de Student para muestras emparejadas y la prueba no paramétrica de rangos con signo de Wilcoxon.

\hfill

Al comparar el modelo ganador, \textttbreak{EncDec-Opt-Small-MT8\_FT\_AR\_Contiguous}, contra los baselines y variantes clave en la Tabla~\ref{tab:resultados-finales-principales}, los resultados son concluyentes:
\begin{itemize}
	\item \textbf{Frente a \CTTransformerBase}: La reducción de error es altamente significativa tanto en la prueba t de Student ($t = -20.55$, $p < 0.0001$) como en la prueba de Wilcoxon ($W = 0.0$, $p \approx 0.00006$).
	\item \textbf{Frente a la mejor variante de encoder (\CTTransformer\ pequeño + Time2Vec + sesgo temporal)}: La ventaja del modelo encoder-decoder es también muy robusta estadísticamente ($t = -16.42$, $p < 0.0001$; Wilcoxon $W = 0.0$, $p \approx 0.00006$).
	\item \textbf{Frente al Modelo lineal temporal}: Se observa significancia estadística rigurosa ($t = -18.28$, $p < 0.0001$; Wilcoxon $W = 0.0$, $p \approx 0.00006$).
	\item \textbf{Frente a STraTS adaptado y Persistencia}: Las diferencias de medias favorecen al modelo ganador con un nivel de confianza superior al 99.9\% ($p < 0.0001$ en la prueba t de Student, y $W = 0.0$, $p \approx 0.00006$ en Wilcoxon).
\end{itemize}

Estos resultados indican que la superioridad del modelo ganador no es producto de fluctuaciones aleatorias en las semillas, sino que existe una diferencia de desempeño sistemática y estadísticamente significativa frente a todas las demás arquitecturas evaluadas bajo los mismos datasets y condiciones de entrenamiento.

\section{Discusión sobre MAPE}
\label{sec:discusion-mape-final}

El MAPE presenta alta variabilidad en las corridas actuales y, en algunas variantes autoregresivas, se registra en columnas auxiliares separadas. Esto es esperable cuando la variable objetivo puede acercarse a cero, porque pequeños denominadores inflan artificialmente el error porcentual. Por ello, el análisis principal se basa en MSE, RMSE y MAE.

\section{Síntesis del capítulo}
\label{sec:sintesis-resultados-final}

Los resultados finales permiten resumir la evaluación en cinco puntos:

\begin{enumerate}
	\item El mejor modelo observado es \textttbreak{EncDec-Opt-Small-MT8\_FT\_AR\_Contiguous}: encoder--decoder pequeño, entrenado con ocho objetivos y ajustado autoregresivamente sobre objetivos contiguos.
	\item Este modelo gana en los 5 datasets y obtiene MSE 0.1979, RMSE 0.4417 y MAE 0.2193.
	\item Dentro de los encoders puros, las variantes pequeñas de \CTTransformer\ son las más competitivas; la mejor por MSE es la versión pequeña con Time2Vec y sesgo temporal, aunque queda prácticamente empatada con \CTTransformerSmall.
	\item Las ablaciones indican que el identificador de token objetivo aporta una señal útil; el efecto aislado del encoding temporal continuo no queda demostrado en este subconjunto.
	\item Aumentar capacidad no mejora el MSE en esta evaluación. Los modelos pequeños entregan la mejor relación entre error, parámetros y tiempo.
\end{enumerate}
